\documentclass[../../main.tex]{subfiles}
\graphicspath{{\subfix{../../image/}}}

\begin{document}

\section{参考文献}\label{section:ch03-references}

\S3.1：本材料的一个好来源是 Courant--Hilbert \cite{C-H}（第 2 章）。

\S3.2：特征微分方程的这一推导见于 Carath\'eodory \cite{C}。\cref{section:ch03-characteristics} 特征定理的证明沿用 Garabedian \cite{G}（第 2 章）、John \cite{J2}（第 1 章）等。Chester \cite{Ch} 与 Sneddon \cite{Sn} 也是一阶 PDE 的好教材。\cref{subsection:ch03-char-examples}的例取自 Zwillinger \cite{Zw}。

\S3.3：关于 Hamilton--Jacobi PDE 的更多内容见 Lions \cite{Li}、Rund \cite{Ru} 与 Benton \cite{Be}。唯一性证明归功于 A. Douglis，取自 \cite{Be}。

\S3.4：见 Lax \cite{Lx1} 与 Smoller \cite{S}（第 15、16 章；定理\cref{theorem:ch03-entropy-uniqueness} 的证明归功于 O. Oleinik，我取自该书）。定理\cref{theorem:ch03-Linf-decay} 与定理\cref{theorem:ch03-L1-decay} 来自 DiPerna（Indiana U. Math. J. 24 (1975), 1047--1071），并感谢 M. Struwe 与 E. Wiedemann 对证明的帮助。关于非线性波的一个好的综合参考文献是 Whitham \cite{Wh}。

\S3.5：习题\cref{exercise:ch03-3} 基于 Miller \cite{Mi}。

\end{document}